\section{应用系统工具：从模型到产品}
\label{sec:tools}

\begin{guide}
\begin{minipage}[c]{0.66\linewidth}
\noindent 只有模型还算不上产品：它要部署运行，要连接外部知识，要嵌入业务流程，有时还要针对特定任务微调。每个环节都从零搭建显然不现实。本节按“负责哪一段链路”把常用工具分为部署工具、微调工具、向量数据库、应用开发框架和编排平台五类，它们各司其职、组合使用。选型时，先确定需求落在哪一段链路，再在同类工具之间比较。
\end{minipage}\hfill
\begin{minipage}[c]{0.3\linewidth}\centering\locator{1}\end{minipage}
\end{guide}

\subsection{一个大模型应用由哪些部分组成}

下面以一个企业内部的知识问答助手为例。用户在网页上提问后，系统需要依次完成以下步骤：判断用户身份和权限，从内部文档中检索相关资料，把问题和资料组织成提示词并交给模型生成回答，把带有引用出处的回答返回给用户，并记录日志以便评估和改进。其中，“运行模型”只是一个环节。

图\,\ref{fig:app-arch} 给出了这类应用的典型分层：\emph{接口层}负责与用户交互；\emph{编排层}负责流程控制、权限校验、工具调用与安全审查；\emph{知识层}提供向量数据库、搜索引擎和文档来源；\emph{模型服务层}由推理引擎运行模型并提供 API；\emph{基础设施层}即第\,\ref{sec:hw}\,节介绍的算力资源。这一分层是对常见组织方式的归纳，并非所有产品都必须恰好分为五层。

% 图 4-6 第二稿：遵循 nature-figure 矢量示意图规范
%  - 一种字族、两种字号：层名 \small 加粗（唯一加粗处），说明与工具 \footnotesize 常规
%  - 全部文字同一深色；唯一强调色只标“工具类别”这一本图要传达的信息
%  - 同一角色的框等宽、上下边对齐；每个框内部顺序相同：层名 → 一行说明 → （底部）工具类别
%  - 连线标签与其他文字保持间距；虚线只表示一种含义：离线
\begin{figure}[htbp]
\centering
\begin{tikzpicture}[fig, x=1mm, y=1mm]
  % \layerbox{x0}{x1}{ytop}{h}{填充}{边框样式}{层名}{说明}{工具类别，可空}
  \newcommand{\layerbox}[9]{%
    \begin{scope}
      \path[clip, rounded corners=1.5pt] (#1,#3) rectangle (#2,#3-#4);
      \fill[#5] (#1,#3) rectangle (#2,#3-#4);
      \if\relax\detokenize{#9}\relax\else
        \fill[accentpale] (#1,#3-#4) rectangle (#2,#3-#4+5);
        \draw[accent, line width=0.4pt] (#1,#3-#4+5) -- (#2,#3-#4+5);
      \fi
    \end{scope}
    \draw[figline, rounded corners=1.5pt, #6] (#1,#3) rectangle (#2,#3-#4);
    \if\relax\detokenize{#9}\relax
      \node[align=center, text=ink] at ({(#1+#2)/2},{#3-#4/2})
        {{\small\bfseries #7}\\[0.5pt]{\footnotesize #8}};
    \else
      \node[align=center, text=ink] at ({(#1+#2)/2},{#3-(#4-5)/2})
        {{\small\bfseries #7}\\[0.5pt]{\footnotesize #8}};
      \node[text=ink, font=\figfont\footnotesize] at ({(#1+#2)/2},{#3-#4+2.5}) {工具：#9};
    \fi
  }
  % ---- 第 1 层 ----
  \layerbox{-55}{55}{0}{11}{figfill}{}{接口层}{网页、移动端、开放 API}{}
  % ---- 第 2 层 ----
  \layerbox{-55}{55}{-19}{16}{figfill}{}{编排层}{工作流与智能体、提示词组织、权限校验、安全审查、日志}{编排平台、应用开发框架}
  % ---- 第 3 层：三个等宽框，等间距 ----
  \layerbox{-55}{-23}{-48}{16}{figfill}{}{知识层}{向量检索、搜索、文档库}{向量数据库}
  \layerbox{-16}{16}{-48}{16}{figfill}{}{模型服务层}{推理引擎、模型 API}{部署工具}
  \layerbox{23}{55}{-48}{16}{white}{dashed}{微调}{离线训练，产出新权重}{微调工具}
  % ---- 第 4 层 ----
  \layerbox{-55}{55}{-76}{11}{white}{}{基础设施层}{计算芯片、存储、互联、驱动与软件栈（第\,4.2\,节）}{}
  % ---- 连线：主路径实线，离线虚线 ----
  \draw[<->, line width=0.6pt] (0,-11) -- (0,-19);
  \draw[<->, line width=0.6pt] (-39,-35) -- (-39,-48) node[midway, right=1.5mm, note] {检索资料};
  \draw[<->, line width=0.6pt] (0,-35)   -- (0,-48)   node[midway, right=1.5mm, note] {生成回答};
  \draw[->, dashed, line width=0.6pt] (23,-56) -- (16,-56);
  \foreach \x in {-39,0,39} \draw[->, line width=0.6pt] (\x,-64) -- (\x,-76);
\end{tikzpicture}
\caption{大模型应用的典型分层与五类工具的位置。每个框底部的色带标出该层对应的工具类别。编排层发起检索与模型调用；微调（虚线框）在线下进行，产出的权重交给模型服务层加载}
\label{fig:app-arch}
\end{figure}


\subsection{五类工具的定位}
\label{subsec:tool-types}

根据所负责的链路，常用工具可分为五类（表\,\ref{tab:tools}）。需要注意，这是按\emph{功能定位}划分的分类，而非严格互斥的技术分类：许多产品同时具备多种能力。例如，某个部署工具既能管理模型，也能对外提供 API；某个编排平台既能搭建知识库，也能集成智能体。选型时应先明确工具负责哪一段链路，再比较具体产品。

\begin{table}[htbp]
\centering
\caption{五类应用系统工具的定位}
\label{tab:tools}
\small
\begin{tabularx}{\linewidth}{@{}l >{\raggedright\arraybackslash}p{30mm} Y Y@{}}
\toprule
\heiti 类型 & \heiti 代表工具 & \heiti 主要功能 & \heiti 容易混淆之处 \\
\midrule
编排平台 & Coze、Dify、FastGPT、MaxKB & 可视化搭建工作流、接入知识库、发布应用 & 不是训练或计算框架 \\
应用开发框架 & LangChain、LlamaIndex、LangGraph & 以代码连接模型、工具与数据，实现控制逻辑 & 不提供算力，也不运行模型 \\
向量数据库 & Milvus、Faiss、Pinecone、Weaviate & 向量存储、近似最近邻检索、元数据过滤 & 只负责检索，不能代替模型生成回答 \\
微调工具 & LLaMA-Factory、Transformers + PEFT、XTuner、Unsloth & 组织数据与训练配置，支持 LoRA 等方法 & 与驱动、编译器不在同一层 \\
部署工具 & Ollama、LM Studio、vLLM、SGLang、TGI & 加载模型、执行推理、提供 API & “方便运行”与“高并发吞吐”是不同目标 \\
\bottomrule
\end{tabularx}
\end{table}

\subsection{部署工具：本地运行与高并发服务}
\label{subsec:deploy}

部署工具可以按目标分为两类。

\paragraph{面向易用的本地运行工具} \term{Ollama} 是开源的跨平台工具，一条命令（如 \texttt{ollama run <模型名>}）即可下载并运行模型，提供兼容 OpenAI 格式的 REST API，并通过量化等手段使模型能够在个人电脑的 CPU 或 GPU 上运行。\term{LM Studio} 则提供图形界面，便于浏览、下载和试用模型。这类工具降低了模型部署运行的门槛，但它们不会改变模型的参数量和知识截止时间，也不会提升模型本身的能力。部分工具支持导入在其他工具中微调好的权重，但这并不等于工具本身完成了微调训练。

\paragraph{面向吞吐的推理引擎} \term{vLLM}、\term{SGLang}、TGI 等推理引擎面向多用户在线服务，其核心是第\,\ref{subsec:infer-opt}\,小节介绍的技术：PagedAttention 管理 KV 缓存，连续批处理调度请求，并集成高效注意力内核与量化支持\cite{kwon2023vllm}。llama.cpp 则以 C/C++ 实现、对 CPU 和消费级硬件友好，常作为本地工具的底层引擎。

二者的取舍可以概括为：面向个人体验、离线开发和小规模试用时，优先考虑易用的本地工具；面向多用户并发、关注吞吐与尾延迟的生产服务时，应选用推理引擎，并按第\,\ref{subsec:metrics}\,小节的方法实测。

\subsection{编排、检索与微调工具}

\paragraph{编排平台与应用开发框架} \term{LangChain} 等应用开发框架以代码形式连接模型、工具和数据源，用“链”或“图”组织多步流程；\term{Coze}、Dify 等低代码平台则通过可视化节点搭建问答、审批和知识库等业务流程。流程设计中有一个重要的权衡\cite{anthropic2024agents}：对于步骤固定、重复执行的任务，采用预先定义的\term{工作流}（\en{workflow}）更便于测试和管控；对于目标多变、步骤难以预先确定的任务，可引入能够自主选择工具、规划步骤的\term{智能体}（\en{agent}），同时必须设置权限边界和停止条件。智能体并非自主性越高越好，自主性越高，出错的代价和外部副作用通常也越大。

\paragraph{向量数据库与检索增强} 在检索增强生成（RAG）应用中，文档先被切分成片段，用嵌入模型转换为向量后存入向量数据库；用户提问时，系统检索与问题最相似的片段，交给模型参考作答。\term{Milvus} 是开源的分布式向量数据库，支持海量向量与元数据的存储、多种索引和相似性检索，采用存储与计算分离的云原生架构\cite{milvus2024}；Faiss 则是一个向量相似性检索\emph{库}，需要嵌入到其他系统中使用，在部署和扩展方式上与完整的数据库不同。需要注意，向量相似并不等于事实正确，实际系统通常还要结合关键词检索、重排序、文档权限控制和回答评估。

\paragraph{微调工具} \term{LLaMA-Factory} 等工具将数据处理、训练配置、LoRA/QLoRA 等参数高效微调方法\cite{hu2021lora,dettmers2023qlora}以及评估导出集成在一起，并提供网页界面，降低了实验管理的门槛。但它们仍然依赖底层框架和硬件，并不会消除优化器状态和激活值带来的开销。要判断一个微调任务能否在单卡上完成，仍需按第\,\ref{subsec:train-mem}\,小节的方法，根据显存、量化方式、训练精度、批大小和序列长度进行估算。

\begin{pitfall}
“先微调还是先做 RAG”取决于要解决的问题。如果痛点是知识需要频繁更新、回答需要给出出处，RAG 通常是成本更低的起点；如果痛点是输出风格、特殊格式或特定任务行为，才进一步考虑微调。二者也可以结合使用。
\end{pitfall}

\begin{keypoints}
\begin{itemize}[leftmargin=1.5em]
  \item “运行模型”只是大模型应用的一个环节，完整的产品还需要接口、编排、知识和基础设施。
  \item 五类工具按负责的链路区分；选型时先确定链路，再比较产品。
  \item 本地运行工具追求易用，推理引擎追求吞吐；工作流追求可控，智能体追求灵活。
\end{itemize}
\end{keypoints}
